\section{Любые два алгебраических замыкания изоморфны}

\begin{theorem}[6.6]
Пусть $\Om$~--- алгебраическое замыкание поля $F$, а $E$~--- алгебраическое расширение $F$. Тогда существует
$F$-гомоморфизм $E\to\Om$; если $E$ тоже алгебраическое замыкание $F$, то каждый такой гомоморфизм~--- изоморфизм.
\end{theorem}

\begin{proof}
\emph{Счётный случай.} Пусть $E=F[\alpha_1,\alpha_2,\dots]$. Продолжим вложение $F\to\Om$ до $F[\alpha_1]$ (пошлём
$\alpha_1$ в любой корень его минимального многочлена в $\Om$), затем до $F[\alpha_1,\alpha_2]$ и т.\,д.: если
$\varphi_n\colon F[\alpha_1,\dots,\alpha_n]\to\Om$ построено и $h$~--- минимальный многочлен $\alpha_{n+1}$ над
$F[\alpha_1,\dots,\alpha_n]$, то у многочлена $\varphi_n(h)$ есть корень в алгебраически замкнутом $\Om$, и по \rf{2.4}
$\varphi_n$ продолжается на $F[\alpha_1,\dots,\alpha_{n+1}]$. Объединение $\varphi_n$~--- искомый гомоморфизм.

\emph{Общий случай.} Пусть $\mathcal S$~--- множество пар $(M,\varphi_M)$, где $F\subset M\subset E$~--- поле, а
$\varphi_M\colon M\to\Om$~--- $F$-гомоморфизм; $(M,\varphi_M)\le(N,\varphi_N)$, если $M\subset N$ и
$\varphi_N|_M=\varphi_M$. Множество $\mathcal S$ непусто (содержит $(F,\text{вложение})$). Если $T$~--- цепь, то
$M'=\bigcup_{M\in T}M$~--- подполе $E$, и правило $\varphi'(x)=\varphi_M(x)$ при $x\in M$ корректно задаёт
гомоморфизм $\varphi'\colon M'\to\Om$; пара $(M',\varphi')$~--- верхняя грань $T$. По лемме Цорна есть максимальный
элемент $(M,\varphi)$. Если $M\ne E$, возьмём $\alpha\in E\setminus M$ и его минимальный многочлен $h$ над $M$; у
$\varphi(h)$ есть корень $\beta\in\Om$, и по \rf{2.4} $\varphi$ продолжается на $M[\alpha]$ правилом $\alpha\mapsto\beta$~---
противоречие с максимальностью. Значит, $M=E$.
\[
\begin{tikzcd}[column sep=large,row sep=large]
F\arrow[r,hook] & M\arrow[r,hook]\arrow[d,"\varphi"'] & M[\alpha]\arrow[dl,dashed,cnew,"{\alpha\,\mapsto\,\beta}"]\arrow[r,hook] & E\\
 & \Om & &
\end{tikzcd}
\]
\figcap{Если $M\ne E$, то $\varphi$ продолжается ещё на один шаг: $\beta$~--- корень многочлена $\varphi(h)$ в
алгебраически замкнутом $\Om$.}

\emph{Изоморфизм.} Пусть $E$ алгебраически замкнуто и $\varphi\colon E\to\Om$~--- $F$-гомоморфизм. Он инъективен, а
поле $\varphi(E)\cong E$ алгебраически замкнуто, так что каждый $f\in F[X]$ раскладывается в $\varphi(E)[X]$ на линейные
множители. Пусть $\alpha\in\Om$ и $f$~--- его минимальный многочлен над $F$. В $\Om[X]$ многочлен $f$ делится на
$X-\alpha$, а его разложение на линейные множители в $\varphi(E)[X]$ является и разложением в $\Om[X]$; по
единственности разложения $\alpha\in\varphi(E)$. Значит, $\varphi$ сюръективен.
\end{proof}

Это типичное применение леммы Цорна: если мы умеем сделать что-то в конечной (или счётной) ситуации, то лемма Цорна
позволяет сделать это в общей.

\begin{remark}[6.7]
В доказательстве \rf{6.6} использована лемма Цорна. Приведём доказательство изоморфности двух алгебраических замыканий
$\Om$ и $\Om'$ поля $F$, использующее лишь принцип ультрафильтра. Для унитарного $f\in F[X]$ обозначим через
$\Om_f\subset\Om$ и $\Om'_f\subset\Om'$ подполя, порождённые над $F$ корнями $f$; это поля разложения $f$ над $F$. Пусть
$H_f$~--- множество $F$-изоморфизмов $\Om_f\to\Om'_f$.

\emph{(a) $H_f$ конечно и непусто.} Непусто по \rf{2.13} (любые два поля разложения изоморфны). Конечно: $F$-гомоморфизм
поля $\Om_f=F[\alpha_1,\dots,\alpha_n]$ определяется образами корней $\alpha_i$, а это корни $f$ в $\Om'$.

\emph{(b) При $f\mid g$ ограничение $H_g\to H_f$, $h\mapsto h|_{\Om_f}$, определено и сюръективно.} Если
$f=\prod_i(X-\alpha_i)$ в $\Om_g[X]$, то, применяя $h$, получаем $f=\prod_i(X-h(\alpha_i))$ в $\Om'[X]$: $h$ переводит
корни $f$ в корни $f$, и $h(\Om_f)=F[h(\alpha_1),\dots,h(\alpha_n)]=\Om'_f$. Сюръективность: пусть $\varphi\in H_f$. Поле
$\Om_g$ порождено над $\Om_f$ корнями $g$, а $g$ раскладывается в $\Om'_g$; присоединяя корни $g$ по одному и каждый
раз применяя \rf{2.4} (образ минимального многочлена очередного корня делит $g$ и потому раскладывается в $\Om'_g$), продолжим
$\varphi$ до гомоморфизма $\hat\varphi\colon\Om_g\to\Om'_g$. Как и выше, $\hat\varphi(\Om_g)=\Om'_g$, так что
$\hat\varphi\in H_g$.

\emph{(c) Компактность.} Снабдим каждое $H_f$ дискретной топологией и положим $H=\prod_fH_f$ (произведение по всем
унитарным $f$). Пространство $H$ компактно (теорема Тихонова для конечных дискретных пространств) и непусто (произведение
непустых конечных множеств); оба факта следуют из принципа ультрафильтра.

\emph{(d)} Для $f\mid g$ множество
\[
H_{f,g}=\{(h_q)_q\in H\mid h_g|_{\Om_f}=h_f\}
\]
замкнуто: это множество точек, где совпадают два непрерывных отображения $H\to H_f$~--- проекция на $H_f$ и проекция на $H_g$ с
последующим ограничением, а $H_f$ хаусдорфово.

\emph{(e) Конечные пересечения непусты.} Пусть даны пары $f_1\mid g_1,\dots,f_n\mid g_n$; положим $k=g_1\cdots g_n$ и
возьмём $\varphi\in H_k$. Возьмём любой элемент $H$ и заменим в нём координаты с номерами $f\mid k$ на
$\varphi|_{\Om_f}$ (это элемент $H_f$ по (b)). Получим $(h_f)\in H$, для которого $h_{g_i}|_{\Om_{f_i}}=
\varphi|_{\Om_{f_i}}=h_{f_i}$, т.\,е.\ $(h_f)\in H_{f_1,g_1}\cap\dots\cap H_{f_n,g_n}$.

\emph{(f) Пересечение всех $H_{f,g}$ непусто.} Иначе дополнения $H\setminus H_{f,g}$ образовали бы открытое покрытие
компактного $H$, из него выделилось бы конечное подпокрытие, и пересечение конечного числа $H_{f,g}$ было бы пусто
вопреки (e).

\emph{(g) Построение изоморфизма.} Пусть $(h_f)\in\bigcap H_{f,g}$, т.\,е.\ $h_g|_{\Om_f}=h_f$ при $f\mid g$. Если
$\alpha\in\Om_f$, то $h_f(\alpha)=h_{\mu_\alpha}(\alpha)$, где $\mu_\alpha$~--- минимальный многочлен $\alpha$ над $F$:
действительно, при $g=f\mu_\alpha$ оба поля $\Om_f$ и $\Om_{\mu_\alpha}$ содержатся в $\Om_g$ и содержат $\alpha$, и
$h_f(\alpha)=h_g(\alpha)=h_{\mu_\alpha}(\alpha)$. Положим $\varphi(\alpha)=h_{\mu_\alpha}(\alpha)$; по только что
доказанному $\varphi(\alpha)=h_f(\alpha)$ для \emph{любого} $f$ с $\alpha\in\Om_f$.
\begin{itemize}[leftmargin=*]
\item $\varphi$~--- гомоморфизм: для $\alpha,\beta\in\Om$ поле $\Om_f$ при $f=\mu_\alpha\mu_\beta$ содержит $\alpha$,
$\beta$, $\alpha+\beta$, $\alpha\beta$, поэтому $\varphi(\alpha+\beta)=h_f(\alpha+\beta)=h_f(\alpha)+h_f(\beta)=
\varphi(\alpha)+\varphi(\beta)$ и так же для произведения; $\varphi(c)=c$ при $c\in F$.
\item $\varphi$ инъективен как гомоморфизм полей.
\item $\varphi$ сюръективен: если $\beta'\in\Om'$ и $\mu$~--- его минимальный многочлен над $F$ (здесь используется,
что $\Om'$ алгебраично над $F$), то $\beta'\in\Om'_\mu=h_\mu(\Om_\mu)$, т.\,е.\ $\beta'=h_\mu(\alpha)$ для некоторого
$\alpha\in\Om_\mu$, и $\varphi(\alpha)=h_\mu(\alpha)=\beta'$.
\end{itemize}
Итак, $\varphi\colon\Om\to\Om'$~--- $F$-изоморфизм.
\end{remark}

\[
\begin{tikzcd}[column sep=large,row sep=small]
 & \Om_f\arrow[r,hook]\arrow[dd,"h_f"',"\cong"] & \Om_g\arrow[r,hook]\arrow[dd,dashed,cnew,"h_g"] & \Om\arrow[dd,dashed,cnew,"\varphi"]\\
F\arrow[ur,hook]\arrow[dr,hook] & & & \\
 & \Om'_f\arrow[r,hook] & \Om'_g\arrow[r,hook] & \Om'
\end{tikzcd}
\]
\figcap{Изоморфизм $h_f$ полей разложения продолжается до $h_g$ при $f\mid g$ (сюръективность $H_g\to H_f$), а
согласованное семейство $(h_f)$, т.\,е.\ точка пересечения всех $H_{f,g}$, склеивается в изоморфизм
$\varphi\colon\Om\to\Om'$.}

Смысл этого доказательства. Унитарные многочлены, упорядоченные по делимости, образуют направленное множество (общее
кратное $f$ и $g$~--- например, $fg$), а $\bigcap H_{f,g}$~--- это обратный предел $\varlim H_f$ конечных непустых
множеств относительно ограничений (ср.~\rf{7.20}--\rf{7.22}). Мы доказали, что такой предел непуст (вариант <<леммы Кёнига>>); для
этого и нужен принцип ультрафильтра. Обратно, каждый $F$-изоморфизм $\Om\to\Om'$ задаёт точку предела
$(\varphi|_{\Om_f})_f$, так что $F$-изоморфизмы $\Om\to\Om'$~--- это в точности точки $\varlim H_f$. Так же устроена
группа Галуа бесконечного расширения (\rf{7.8}, \rf{7.26}).

\begin{remark}[6.8: предостережение]
Даже для конечного поля $F$ изоморфизмов одного алгебраического замыкания на другое несчётно много, и ни один из них
не выделен. Действительно, если $\varphi_0\colon\Om\to\Om'$~--- один $F$-изоморфизм, то все остальные имеют вид
$\varphi_0\circ\sigma$, $\sigma\in\Aut(\Om/F)$, а, например, $\Aut(\Fpb/\F_p)\cong\Zh$ (\rf{7.17}) несчётна (\rf{7.26b}). Поэтому
говорить о <<единственном>> алгебраическом замыкании $F$ неаккуратно: можно лишь сказать, что для любых двух
алгебраических замыканий $\Om$, $\Om'$ поля $F$ существует (в предположении принципа ультрафильтра) $F$-изоморфизм
$\Om\to\Om'$.
\end{remark}

\begin{aside}[6.8a: каноническое алгебраическое замыкание]
Единственность бывает двух уровней.
\begin{itemize}[leftmargin=*]
\item \emph{С точностью до единственного изоморфизма.} Так обстоит дело с объектами, заданными универсальным
свойством: поле частных кольца, фактор-кольцо, тензорное произведение. Любые две реализации отождествляются ровно одним
изоморфизмом, согласованным со структурными отображениями ($A\to\operatorname{Frac}A$, $A\to A/I$,
$V\times W\to V\otimes W$), и можно говорить о <<том самом>> объекте и его элементах.
\item \emph{С точностью до неединственного изоморфизма}~--- случай алгебраического замыкания. Здесь даже
$F$-изоморфизмы, т.\,е.\ согласованные с вложениями $F\to\Om$, не единственны. В абстрактном алгебраическом замыкании
$\Q$ нельзя указать, какой из двух корней $X^2+1$ является <<тем самым $i$>>: корни $i$ и $-i$ не различаются никакими
свойствами над $\Q$ (автоморфизм, переставляющий их, сохраняет все такие свойства). Двух экземпляров алгебраического
замыкания нельзя отождествить без произвольного выбора.
\end{itemize}
Поэтому \emph{каноническим} называют объект, заданный без произвольных выборов. Если $F$ лежит в фиксированном
алгебраически замкнутом поле $\Om$, то алгебраическое замыкание $F$ в $\Om$ (\rf{1.46}, \rf{1.47}) каноническое: это конкретное
подмножество $\Om$. Так, $\Qal\subset\CC$~--- <<настоящее>> алгебраическое замыкание $\Q$, в котором $i$, вещественное
$\sqrt[3]2$ и $e^{2\pi i/5}$~--- конкретные элементы, а комплексное сопряжение~--- конкретный элемент
$\Gal(\Qal/\Q)$. Каноничность здесь относительная: она опирается на то, что $F$ уже лежит в $\CC$. Абстрактное поле
$\Q[X]/(X^3-2)$ вкладывается в $\CC$ тремя способами, а само $\CC$ имеет огромную группу автоморфизмов (её мощность
$2^{2^{\aleph_0}}$; все они, кроме тождественного и комплексного сопряжения, разрывны и строятся с помощью аксиомы
выбора). Следствие для теории Галуа: абсолютная группа
Галуа $\Gal(\Fsep/F)$ (сепарабельное замыкание $\Fsep$~--- см.~\rf{6.9}, \rf{6.10}) определена лишь с точностью до внутреннего автоморфизма (\rf{7.6a}). Без аксиомы выбора дело обстоит
ещё хуже: у $\Q$ могут быть и неизоморфные алгебраические замыкания (\rf{6.4a}).
\end{aside}
